% Options for packages loaded elsewhere
\PassOptionsToPackage{unicode}{hyperref}
\PassOptionsToPackage{hyphens}{url}
%
\documentclass[
]{article}
\usepackage{amsmath,amssymb}
\usepackage{iftex}
\ifPDFTeX
  \usepackage[T5]{fontenc}
  \usepackage[utf8]{inputenc}
  \usepackage{textcomp} % provide euro and other symbols
\else % if luatex or xetex
  \usepackage{unicode-math} % this also loads fontspec
  \defaultfontfeatures{Scale=MatchLowercase}
  \defaultfontfeatures[\rmfamily]{Ligatures=TeX,Scale=1}
\fi
\usepackage{lmodern}
\ifPDFTeX\else
  % xetex/luatex font selection
\fi
% Use upquote if available, for straight quotes in verbatim environments
\IfFileExists{upquote.sty}{\usepackage{upquote}}{}
\IfFileExists{microtype.sty}{% use microtype if available
  \usepackage[]{microtype}
  \UseMicrotypeSet[protrusion]{basicmath} % disable protrusion for tt fonts
}{}
\makeatletter
\@ifundefined{KOMAClassName}{% if non-KOMA class
  \IfFileExists{parskip.sty}{%
    \usepackage{parskip}
  }{% else
    \setlength{\parindent}{0pt}
    \setlength{\parskip}{6pt plus 2pt minus 1pt}}
}{% if KOMA class
  \KOMAoptions{parskip=half}}
\makeatother
\usepackage{xcolor}
\usepackage{graphicx}
\makeatletter
\def\maxwidth{\ifdim\Gin@nat@width>\linewidth\linewidth\else\Gin@nat@width\fi}
\def\maxheight{\ifdim\Gin@nat@height>\textheight\textheight\else\Gin@nat@height\fi}
\makeatother
% Scale images if necessary, so that they will not overflow the page
% margins by default, and it is still possible to overwrite the defaults
% using explicit options in \includegraphics[width, height, ...]{}
\setkeys{Gin}{width=\maxwidth,height=\maxheight,keepaspectratio}
% Set default figure placement to htbp
\makeatletter
\def\fps@figure{htbp}
\makeatother
\ifLuaTeX
  \usepackage{luacolor}
  \usepackage[soul]{lua-ul}
\else
  \usepackage{soul}
\fi
\setlength{\emergencystretch}{3em} % prevent overfull lines
\providecommand{\tightlist}{%
  \setlength{\itemsep}{0pt}\setlength{\parskip}{0pt}}
\setcounter{secnumdepth}{-\maxdimen} % remove section numbering
\ifLuaTeX
  \usepackage{selnolig}  % disable illegal ligatures
\fi
\usepackage{bookmark}
\IfFileExists{xurl.sty}{\usepackage{xurl}}{} % add URL line breaks if available
\urlstyle{same}
\hypersetup{
  hidelinks,
  pdfcreator={LaTeX via pandoc}}

\author{}
\date{}

\usepackage[vietnamese]{babel}
\usepackage{ulem}
\let\ul\uline
\begin{document}
\ul{Buổi 03 -- Ngày 02-08-2026 -- môn Đại số tuyến tính -- lớp
MA003.F32.LT.CNTT}

\textbf{6/ \ul{MA TRẬN VUÔNG KHẢ NGHỊCH}}

a/ Cho ma trận $A$
vuông cấp $n$ trên
\emph{F}. Ta nói $A$
khả nghịch (\emph{invertible matrix}) nếu có ma trận vuông
$B$ (cấp
$n$) sao cho

$AB = BA = I_n$

b/ Ma trận $B$ (nếu
có) thì duy nhất, và ta ký hiệu là
$A^{-1}$ và ta gọi
$A^{-1}$ là ma trận
nghịch đảo (\emph{inverse matrix}) của
$A$.

c/ Khi $A$ khả
nghịch, ta có thêm lũy thừa nguyên âm của
$A$ như sau (bên
cạnh lũy thừa nguyên, không âm đã có trước đó)

$A^{-m} = (A^{-1})^m$ là nghịch đảo
của ma trận $A$;

$A^0 = I_n$

$A^{-n} = (A^n)^{-1}$

d/ \emph{\ul{Cách nhận diện ma trận khả nghịch và tìm ma trận nghịch đảo
(nếu có)}}:

Cho ma trận $A$
vuông cấp $n$ trên
\emph{F}. Ta kiểm tra xem
$A$ có khả nghịch
không và tìm ma trận nghịch đảo (nếu có) như sau:

Ta viết dạng ma trận hóa:
$(A|I_n)$ rồi dùng các
phép biến đổi sơ cấp trên dòng thích hợp để xây dựng trong
$A$ các cột chuẩn
theo thứ tự từ trái \textbf{$\rightarrow$} phải (tương tự như quá trình giải hệ pt
tuyến tính bằng pp GAUSS-JORDAN). Cụ thể như sau:

$(A|I_n) \xrightarrow{b_i} (A_1|B_1) \xrightarrow{b_2} \dots \xrightarrow{b_k} (I_n|A^{-1})$ (cặp ma trận
sau cùng là $A_k$ và
$B_k$ sau khi đã xét
xong cột cuối cùng của
$A$).

\emph{\ul{TH1}}: Nếu
$A_k \neq I_n$ , ta nói
$A$ không khả
nghịch và không có ma trận nghịch đảo
$A^{-1}$

\emph{\ul{TH2}}: Nếu
$A_k = I_n$, ta nói
$A$ khả nghịch và
$A^{-1} = B_1$.

\ul{Ví dụ mẫu}: Cho ma trận thực
$$A = \begin{pmatrix} -3 & 4 & 6 \\ 0 & -1 & -1 \\ 2 & -3 & -4 \end{pmatrix}$$. Hỏi
$A$ có khả nghịch
không? Vì sao? Nếu có, tìm
$A^{-1}, A^{-2}$ và
$A^{-4}$.

\ul{Giải}: Ta có

$$(A \mid I_3) = \begin{pmatrix} -3 & 4 & 6 & \mid & 1 & 0 & 0 \\ 0 & -1 & -1 & \mid & 0 & 1 & 0 \\ 2 & -3 & -4 & \mid & 0 & 0 & 1 \end{pmatrix} \xrightarrow[(3)\rightarrow(3)-2(1)\text{mới}]{(1)\rightarrow-1(1)-(3)} \begin{pmatrix} 1 & -1 & -2 & \mid & -1 & 0 & -1 \\ 0 & -1 & -1 & \mid & 0 & 1 & 0 \\ 0 & -1 & 0 & \mid & 2 & 0 & 3 \end{pmatrix}$$

$$\xrightarrow[(2)\rightarrow-1(2)]{\substack{(3)\rightarrow(3)-(2) \\ (1)\rightarrow(1)-(2)}} \begin{pmatrix} 1 & 0 & -1 & \mid & -1 & -1 & -1 \\ 0 & 1 & 1 & \mid & 0 & -1 & 0 \\ 0 & 0 & 1 & \mid & 2 & -1 & 3 \end{pmatrix} \xrightarrow[(1)\rightarrow(1)+(3)]{(2)\rightarrow(2)-(3)} \begin{pmatrix} 1 & 0 & 0 & \mid & 1 & -2 & 2 \\ 0 & 1 & 0 & \mid & -2 & 0 & -3 \\ 0 & 0 & 1 & \mid & 2 & -1 & 3 \end{pmatrix}$$

Do ma trận $A_k = I_n$ nên
$A$ khả nghịch, và
ta có $$A^{-1} = \begin{pmatrix} 1 & -2 & 2 \\ -2 & 0 & -3 \\ 2 & -1 & 3 \end{pmatrix}$$, và

$$A^{-2} = A^{-1} . A^{-1} = \begin{pmatrix} 1 & -2 & 2 \\ -2 & 0 & -3 \\ 2 & -1 & 3 \end{pmatrix} \begin{pmatrix} 1 & -2 & 2 \\ -2 & 0 & -3 \\ 2 & -1 & 3 \end{pmatrix} = \begin{pmatrix} 9 & -4 & 14 \\ -8 & 7 & -13 \\ 10 & -7 & 16 \end{pmatrix}$$

$A^{-4} = A^{-2} . A^{-2} = \dots$

\emph{\ul{Lưu ý}}:

\begin{itemize}
\item
  Nếu ma trận $A$
  vuông cấp $n$ có
  tính chất $A^m = O$,
  với $m$ là một số
  nguyên dương (nào đó) thì ta nói
  $A$ là ma trận
  \emph{lũy linh}.
\item
  Nếu ma trận $A$
  vuông cấp $n$ có
  tính chất $A^2 = A$ thì
  ta gọi $A$ là ma
  trận \emph{lũy đẳng.}
\item
  Nếu ma trận $A$
  vuông cấp $n$ có
  tính chất $A^{-1} = A^T$ thì
  ta gọi $A$ là ma
  trận \emph{trực giao.}
\end{itemize}

Trong đó $a_{ij} = -a_{ji}$ (T =
Transposed matrix), ta gọi
$A^T$ là ma trận
chuyển vị của $A$.

Ta lấy dòng viết thành cột, cột viết thành dòng

\ul{Ví dụ}: Ta có

$$A = \begin{pmatrix} -3 & 2 & 5 & 4 \\ 9 & -7 & 1 & 6 \\ 5 & 3 & 0 & -8 \end{pmatrix} \Rightarrow A^T = \begin{pmatrix} -3 & 9 & 5 \\ 2 & -7 & 3 \\ 5 & 1 & 0 \\ 4 & 6 & -8 \end{pmatrix}$$

Cách bấm MT bỏ túi tìm
$A^{-1}, A^T$.

(minh họa bằng dòng máy tính CASIO fx-570ES PLUS)

MODE \textbf{$\rightarrow$} 6: MATRIX

Matrix?

1: MatA 2: MatB

3: MatC

\begin{itemize}
\item
  Chọn 1: MatA
\item
  1: 3x3
\item
  Nhập ma trận $A$
  (bấm dấu = sau khi di chuyển chuột đến vị trí tương ứng và nhập số
  vào)
\item
  Nhập xong bấm AC / (dòng máy CASIO 580 \textbf{$\rightarrow$} bấm OPTN)
\item
  Bấm SHIFT + 4 \textbf{$\rightarrow$} 3: MatA \textbf{$\rightarrow$} Bấm
  ``$X^{-1}$''
  \textbf{$\rightarrow$} bấm ``='' cho kết quả
  $A^{-1}$
\item
  Lấy kết quả $A^{-1} * A^{-1}$
  hoặc $(A^{-1})^T$ sau đó
  bấm ``='' cho kết quả
  $A^{-2}$.
\end{itemize}

\emph{\ul{Bài tập tương tự}}:

\ul{Bài 1}: Cho ma trận thực
$$B = \begin{pmatrix} 5 & 8 & 9 \\ -12 & -7 & 12 \\ -1 & 5 & 4 \end{pmatrix}$$. Hỏi
$A$ có khả nghịch
không? Vì sao? Nếu có, tìm
$A^3, A^{-1}, A^{-2}, A^{-3}$.

\ul{Bài 2}: (yêu cầu như bài 1), với
$$A = \begin{pmatrix} -2 & 7 & 1 \\ 22 & -71 & -12 \\ -9 & 31 & 5 \end{pmatrix}$$

\ul{Bài 3}: (yêu cầu như bài 1), với
$$B = \begin{pmatrix} 13 & 8 & -12 \\ -12 & -7 & 12 \\ 6 & 4 & -5 \end{pmatrix}$$

\ul{Bài 4}: (yêu cầu như bài 1), với
$$A = \begin{pmatrix} -1 & 4 & -3 \\ -2 & 1 & 2 \\ 5 & -3 & 6 \end{pmatrix}$$

\ul{Bài 5}: (yêu cầu như bài 1), với
$$A = \begin{pmatrix} 1 & 0 & 2 \\ 2 & -1 & 3 \\ -3 & 2 & -5 \end{pmatrix}$$

\ul{Bài 6}: (yêu cầu như bài 1), với
$$A = \begin{pmatrix} 2 & 7 & -8 \\ 1 & 4 & -5 \\ -3 & -9 & 8 \end{pmatrix}$$

\ul{Bài 7} (yêu cầu như bài 1), với
$$A = \begin{pmatrix} 15 & -10 & 5 \\ -7 & 8 & -2 \\ -5 & 3 & -1 \end{pmatrix}$$

\emph{\textbf{\ul{Ứng dụng: giải phương trình ma trận}}}

\textbf{a/ \ul{Phương trình AX = B} (*)}

Trong đó: $A$ là ma
trận vuông cấp $m$

$B$ là ma trận cấp
$m \times n$

$X$ là ma trận ẩn
cần tìm, cấp $m \times n$

Nếu
$A$ không khả
nghịch, thì không tồn tại
$A^{-1}$ và phương trình
vô nghiệm

Nếu $A$ khả nghịch,
ta tìm $A^{-1}$, và pt
có nghiệm duy nhất là:

$X = A^{-1}B$

\ul{Ví dụ 1}: Giải pt ma trận sau trên
$\mathbb{R}$, với
$X$ là ma trận ẩn
cần tìm:

$$\begin{pmatrix} -8 & -3 \\ 5 & 2 \end{pmatrix} X = \begin{pmatrix} -8 & 7 & 10 & -9 \\ 6 & 2 & -4 & 5 \end{pmatrix}$$

b/ \textbf{\ul{Phương trình XA = B} (**)}

Trong đó: $A$ là ma
trận vuông cấp $n$

$B$ là ma trận cấp
$m \times n$

$X$ là ma trận ẩn
cần tìm, cấp $m \times n$

Nếu $A$ không khả
nghịch, thì không tồn tại
$\exists A^{-1}$ và phương trình
vô nghiệm

Nếu $A$ khả nghịch,
ta tìm $A^{-1}$, và pt
có nghiệm duy nhất là:

$X = BA^{-1}$

\ul{Ví dụ 2}: Giải pt ma trận sau trên
$\mathbb{R}$, với
$X$ là ma trận ẩn
cần tìm:

$$X \begin{pmatrix} -8 & -5 \\ 5 & 3 \end{pmatrix} = \begin{pmatrix} -6 & 10 \\ 8 & 3 \\ 5 & -7 \\ 2 & -1 \end{pmatrix}$$

c/ \textbf{\ul{Phương trình AXC = B} (***)}

Trong đó: $A$ là ma
trận vuông cấp $m$

$C$ là ma trận
vuông cấp $n$

$B$ là ma trận cấp
$m \times n$

$X$ là ma trận ẩn
cần tìm, cấp $m \times n$

Nếu $A$ không khả
nghịch, thì không tồn tại
$\exists A^{-1}$ và phương trình
vô nghiệm

Hay $C$ không khả
nghịch, thì không tồn tại
$\exists C^{-1}$ và phương trình
vô nghiệm

Nếu $A$ khả nghịch,
ta tìm $A^{-1}$, và
$C$ khả nghịch, ta
tìm $C^{-1}$ pt có
nghiệm duy nhất là:

$X = A^{-1} B C^{-1}$

\ul{Ví dụ 3}: Giải pt ma trận sau trên
$\mathbb{R}$, với
$X$ là ma trận ẩn
cần tìm:

$$\begin{pmatrix} -7 & -3 \\ 9 & 4 \end{pmatrix} X \begin{pmatrix} -17 & -25 \\ 2 & 3 \end{pmatrix} = \begin{pmatrix} -8 & 15 \\ 6 & -2 \end{pmatrix}$$

\emph{\ul{Lưu ý}}: Khi
$A$ là ma trận
vuông cấp 2, ta xét
$$A = \begin{pmatrix} a & b \\ c & d \end{pmatrix}$$, ta sẽ kiểm
chứng xem $A$ có
khả nghịch không? Và tìm ma trận nghịch đảo
$A^{-1}$ (nếu có) như
sau:

Ta tìm $k = ad - bc$.

Nếu $k = 0$ thì
$A$ không khả
nghịch và không tồn tại
$A^{-1}$.

Nếu $k \neq 0$ thì
$A$ khả nghịch và
$$A^{-1} = \frac{1}{k} \begin{pmatrix} d & -b \\ -c & a \end{pmatrix}$$.

\ul{Ví dụ}: Cho $$A = \begin{pmatrix} -7 & -3 \\ 9 & 4 \end{pmatrix} \Rightarrow k = 7.4 - (-3).9 = -1 \neq 0 \Rightarrow \exists A^{-1}$$
khả nghịch,

và $$A^{-1} = \frac{1}{k} \begin{pmatrix} d & -b \\ -c & a \end{pmatrix} = \frac{1}{-1} \begin{pmatrix} 4 & 3 \\ -9 & -7 \end{pmatrix} = \begin{pmatrix} -4 & -3 \\ 9 & 7 \end{pmatrix}$$

Cho $$C = \begin{pmatrix} -17 & -25 \\ 2 & 3 \end{pmatrix} \Rightarrow k = (-17).3 - (-25).2 = -1 \neq 0 \Rightarrow \exists C^{-1}$$ khả nghịch,

và $$C^{-1} = \frac{1}{k} \begin{pmatrix} d & -b \\ -c & a \end{pmatrix} = \frac{1}{-1} \begin{pmatrix} 3 & 25 \\ -2 & -17 \end{pmatrix} = \begin{pmatrix} -3 & -25 \\ 2 & 17 \end{pmatrix}$$

Cho $$C = \begin{pmatrix} -2 & -3 \\ 7 & 11 \end{pmatrix} \Rightarrow k = (-2).11 - (-3).7 = -1 \neq 0 \Rightarrow \exists C^{-1}$$ không khả
nghịch và không tồn tại
$C^{-1}$.

\ul{Ví dụ 4}: Giải phương trình ma trận sau trên
$\mathbb{R}$, với
$X$ là ma trận ẩn
cần tìm

\begin{quote}
$$\begin{pmatrix} -3 & -4 & -9 \\ 2 & 1 & 2 \\ -7 & 1 & 4 \end{pmatrix} X \begin{pmatrix} -2 & -3 \\ 7 & 11 \end{pmatrix} = \begin{pmatrix} -10 & 8 \\ 6 & -5 \\ 4 & 9 \end{pmatrix}$$
\end{quote}

\ul{Ví dụ 5}: (yêu cầu như ví dụ 4)

$$\begin{pmatrix} -1 & 4 & -3 \\ -1 & 3 & -2 \\ 1 & -1 & 1 \end{pmatrix} X \begin{pmatrix} 1 & 2 & 1 \\ 1 & 3 & 3 \\ 1 & 3 & 4 \end{pmatrix} = \begin{pmatrix} -5 & 8 & -7 \\ 4 & -1 & 6 \\ 9 & 10 & -3 \end{pmatrix}$$

\ul{Ví dụ 6}: (yêu cầu như ví dụ 4)

$$\begin{pmatrix} 3 & -5 & 3 \\ -1 & 3 & -2 \\ 0 & -1 & 1 \end{pmatrix} X \begin{pmatrix} 2 & 7 & -8 \\ 1 & 4 & -5 \\ -3 & -9 & 8 \end{pmatrix} = \begin{pmatrix} 6 & -2 & 10 \\ -3 & 5 & -4 \\ -7 & 1 & -6 \end{pmatrix}$$

\ul{Giải ví dụ mẫu}

Giải pt ma trận:

$$\begin{pmatrix} 4 & 3 \\ 3 & 2 \end{pmatrix} X \begin{pmatrix} 7 & 5 \\ 3 & 2 \end{pmatrix} = \begin{pmatrix} 1 & 2 \\ -1 & 0 \end{pmatrix}$$

Ta có $$A = \begin{pmatrix} 4 & 3 \\ 3 & 2 \end{pmatrix} \Rightarrow k = 4.2 - 3.3 = -1 \neq 0 \Rightarrow \exists A^{-1}$$ khả
nghịch và

$$A^{-1} = \frac{1}{k} \begin{pmatrix} d & -b \\ -c & a \end{pmatrix} = \frac{1}{-1} \begin{pmatrix} 2 & -3 \\ -3 & 4 \end{pmatrix} = \begin{pmatrix} -2 & 3 \\ 3 & -4 \end{pmatrix}$$

Và $$C = \begin{pmatrix} 7 & 5 \\ 3 & 2 \end{pmatrix} \Rightarrow k = 7.2 - 3.5 = -1 \neq 0 \Rightarrow \exists C^{-1}$$ khả nghịch
và

$$C^{-1} = \frac{1}{k} \begin{pmatrix} d & -b \\ -c & a \end{pmatrix} = \frac{1}{-1} \begin{pmatrix} 2 & -5 \\ -3 & 7 \end{pmatrix} = \begin{pmatrix} -2 & 5 \\ 3 & -7 \end{pmatrix}$$

Do \emph{A} và \emph{C} đều khả nghịch nên pt có nghiệm duy nhất là:

$$= \begin{pmatrix} -5 & -4 \\ 7 & 6 \end{pmatrix} \begin{pmatrix} -2 & 5 \\ 3 & -7 \end{pmatrix} = \begin{pmatrix} -2 & 3 \\ 4 & -7 \end{pmatrix}$$

$$ = \begin{pmatrix} -5 & -4 \\ 7 & 6 \end{pmatrix} \begin{pmatrix} -2 & 5 \\ 3 & -7 \end{pmatrix} = \begin{pmatrix} -2 & 3 \\ 4 & -7 \end{pmatrix} $$.

\ul{Ví dụ 7}: Giải phương trình ma trận sau trên
$\mathbb{R}$, với
$X$ là ma trận ẩn
cần tìm

$X^2 = I_2$

\ul{Gợi ý}: Đặt
$$X = \begin{pmatrix} a & b \\ c & d \end{pmatrix}$$ nên pt tương
đương

$$\begin{pmatrix} a & b \\ c & d \end{pmatrix} \cdot \begin{pmatrix} a & b \\ c & d \end{pmatrix} = \begin{pmatrix} 1 & 0 \\ 0 & 1 \end{pmatrix} \Leftrightarrow \begin{pmatrix} a^2 + bc & ab + bd \\ ac + cd & bc + d^2 \end{pmatrix} = \begin{pmatrix} 1 & 0 \\ 0 & 1 \end{pmatrix}$$

Đồng nhất thức 2 vế, ta có hệ pt

$$\begin{cases} a^2 + bc = 1 \\ b(a + d) = 0 \\ c(a + d) = 0 \\ bc + d^2 = 1 \end{cases}$$

(SV tự làm tiếp nhé\ldots)

\ul{Ví dụ 8}: Giải phương trình ma trận sau trên
$\mathbb{R}$, với
$X$ là ma trận ẩn
cần tìm

$$\begin{pmatrix} 1 & 1 & 1 & 1 \\ 1 & -1 & 1 & -1 \\ 1 & -1 & -1 & 1 \\ 1 & 1 & -1 & -1 \end{pmatrix} X \begin{pmatrix} 13 & -16 & 3 \\ -7 & 8 & -2 \\ -3 & 3 & -1 \end{pmatrix} = \begin{pmatrix} -3 & 2 & -5 \\ 8 & -4 & 6 \\ 1 & 10 & -7 \\ 3 & -8 & 9 \end{pmatrix}$$

\ul{Ví dụ 9}: Giải phương trình ma trận sau trên
$\mathbb{R}$, với
$X$ là ma trận ẩn
cần tìm

$$\begin{pmatrix} 2 & 1 \\ 1 & 2 \end{pmatrix} X - X \begin{pmatrix} 1 & -1 \\ -1 & 1 \end{pmatrix} = \begin{pmatrix} 1 & 1 \\ 1 & -1 \end{pmatrix}$$

\ul{Gợi ý}: $X$ có
2 dòng $X$ có 2
cột, nên suy ra
$$X = \begin{pmatrix} a & b \\ c & d \end{pmatrix}$$

Ta có pt tương đương

$$\begin{pmatrix} 2 & 1 \\ 1 & 2 \end{pmatrix} \begin{pmatrix} a & b \\ c & d \end{pmatrix} - \begin{pmatrix} a & b \\ c & d \end{pmatrix} \begin{pmatrix} 1 & -1 \\ -1 & 1 \end{pmatrix} = \begin{pmatrix} 1 & 1 \\ 1 & -1 \end{pmatrix}$$

$$\Leftrightarrow \begin{pmatrix} 2a+c & 2b+d \\ a+2c & b+2d \end{pmatrix} - \begin{pmatrix} a-b & -a+b \\ c-d & -c+d \end{pmatrix} = \begin{pmatrix} 1 & 1 \\ 1 & -1 \end{pmatrix}$$

Viết hệ dưới dạng ma trận hóa ta được:

$$\Leftrightarrow \begin{pmatrix} a+b+c & a+b+d \\ a+c+d & b+c+d \end{pmatrix} = \begin{pmatrix} 1 & 1 \\ 1 & -1 \end{pmatrix}$$

(các bạn SV tự làm tiếp nhé )

\ul{Ví dụ 10}: Giải phương trình ma trận sau trên
$$\Leftrightarrow \begin{cases} a+b+c = 1 \\ a+b+d = 1 \\ a+c+d = 1 \\ b+c+d = -1 \end{cases}$$, với
$$\begin{pmatrix} 1 & 1 & 1 & 0 & 1 \\ 1 & 1 & 0 & 1 & 1 \\ 1 & 0 & 1 & 1 & 1 \\ 0 & 1 & 1 & 1 & -1 \end{pmatrix}$$ là ma trận ẩn
cần tìm

$\mathbb{R}$

\textbf{* \ul{ĐỊNH THỨC MA TRẬN VUÔNG}}

Cho ma trận $X$
vuông cấp $$\begin{pmatrix} 3 & -1 \\ 1 & 3 \end{pmatrix} X + X \begin{pmatrix} -2 & 1 \\ 0 & -1 \end{pmatrix} = \begin{pmatrix} 5 & -8 \\ 3 & -4 \end{pmatrix}$$,
trên \emph{F}. Ta luôn có một con số
$A$ đặc trưng duy
nhất cho $n$, và
ta gọi là định thức của
$X$ (determinant
of $X$), và kí
hiệu là:

$c_A$

\emph{\ul{Cách tính}}

\ul{TH1}: $X$ là
ma trận vuông cấp 1, nghĩa là
$A$

Ví dụ:

\begin{quote}
$c_A = \det(A) = |A|$
\end{quote}

$A = (-7) \Rightarrow \det(A) = |A| = -7$

\ul{TH2}: $X$ là
ma trận vuông cấp 2, nghĩa là

\begin{quote}
$$B = \begin{pmatrix} 3 \\ 5 \end{pmatrix} \Rightarrow \det(B) = |B| = \frac{3}{5}$$
\end{quote}

Ví dụ:

\begin{quote}
$$B = \begin{pmatrix} 3 \\ 5 \end{pmatrix} \Rightarrow \det(B) = |B| = \frac{3}{5}$$
\end{quote}

$$A = \begin{pmatrix} a & b \\ c & d \end{pmatrix} \Rightarrow \det(A) = |A| = ad - bc$$

\ul{TH3}: $X$ là
ma trận vuông cấp 3, nghĩa là
$$A = \begin{pmatrix} -8 & 3 \\ 10 & -2 \end{pmatrix} \Rightarrow \det(A) = |A| = (-8).(-2) - 10.3 = -14$$, ta có quy tắc
SARRUS (quy tắc 6 đường chéo) như sau:

Xét $$B = \begin{pmatrix} -2 & 3 \\ 5 & 7 \end{pmatrix} \Rightarrow \det(B) = |B| = (-2).7 - 5.3 = -29$$

Ta có

$\det(A) = |A| = [a.q.w + b.r.u + c.p.v] - [c.q.u + a.r.v + b.p.w]$

\ul{Ví dụ mẫu}: Cho
$\det(A) = |A| = [a.q.w + b.r.u + c.p.v] - [c.q.u + a.r.v + b.p.w]$

$\Rightarrow \det(A) = |A| = [(-3).6.8 + 5.(-1).4 + (-2).10.(-5)] - [(-2).6.4 + 5.10.8 + (-3).(-1).(-5)] = -401$

(bấm MT bỏ túi: tìm
$\Rightarrow \det(A) = |A| = [(-3).6.8 + 5.(-1).4 + (-2).10.(-5)] - [(-2).6.4 + 5.10.8 + (-3).(-1).(-5)] = -401$) bằng dòng máy
CASIO fx570ES PLUS

MODE $\rightarrow$ 6: MATRIX $\rightarrow$ Màn hình xuất hiện

Matrix?

1: MatA 2: MatB

3: MatC

\begin{itemize}
\item
  Chọn 1 $\rightarrow$ chọn 1:3x3 $\rightarrow$ Nhập ma trận
  $\det(A) = |A|$ (di chuyển
  con trỏ chuột đến từng vị trí ma trận, nhập số rồi bấm dấu =) $\rightarrow$ Nhập
  xong, bấm AC $\rightarrow$ SHIFT + 4 $\rightarrow$ 7: det $\rightarrow$ det(\ldots$\rightarrow$ SHIFT + 4 $\rightarrow$ 3: MatA $\rightarrow$
  ``)'' $\rightarrow$ bấm dấu ``='' $\rightarrow$ cho ra kết quả là =
  $A$.
\end{itemize}

\ul{Bài tập tương tự}:

Tính định thức các ma trận sau:

\ul{Bài 1}: $-401$

\ul{Bài 2}: $$A = \begin{pmatrix} 4 & -7 & 2 \\ 5 & -8 & -1 \\ 6 & 9 & 3 \end{pmatrix}$$

\ul{Bài 3}: $$A = \begin{pmatrix} -2 & 7 & -5 \\ 8 & 3 & 1 \\ 4 & -6 & 9 \end{pmatrix}$$

\ul{Bài 4}: $$A = \begin{pmatrix} -6 & 5 & -1 \\ 8 & -3 & -2 \\ -4 & -9 & 7 \end{pmatrix}$$

\ul{Bài 5}: $$A = \begin{pmatrix} 3 & -4 & -7 \\ -8 & 9 & -5 \\ -1 & -6 & 2 \end{pmatrix}$$

\ul{TH4}: $$A = \begin{pmatrix} 9 & -2 & 3 \\ 8 & 1 & -7 \\ -4 & -6 & 5 \end{pmatrix}$$ là
ma trận vuông cấp
$$A = \begin{pmatrix} 7 & -3 & -8 \\ -2 & 9 & 1 \\ -5 & 4 & -6 \end{pmatrix}$$, với
$n$

\ul{Ký hiệu}:

$n \ge 3$ma trận
$$A = \begin{pmatrix} 9 & -2 & 3 \\ 8 & 1 & -7 \\ -4 & -6 & 5 \end{pmatrix}$$ xóa bỏ đi dòng
$A(i,j) =$ và cột
$(j)$

$c_{ij}^A = (-1)^{i+j} \det[A(i|j)]$, với
$1 \le i, j \le n$

\ul{Ví dụ mẫu}: Cho
$$A = \begin{pmatrix} -7 & 6 & 3 \\ 1 & 4 & -2 \\ 9 & -5 & 10 \end{pmatrix}$$. Tìm
$$\begin{cases} A(2|3) = ? \Rightarrow c_{23}^A = ? \\ A(3|1) = ? \Rightarrow c_{31}^A = ? \\ A(1|2) = ? \Rightarrow c_{12}^A = ? \end{cases}$$

\ul{Giải}:

Ta có:

\begin{quote}
$$A(2|3) = \begin{pmatrix} -5 & 1 \\ 9 & 10 \end{pmatrix} \Rightarrow c_{23}^A = (-1)^{2+3} \det[A(2|3)] = - \begin{vmatrix} -5 & 1 \\ 9 & 10 \end{vmatrix} = -(-50 - 9) = 59$$
\end{quote}

$$A(3|1) = \begin{pmatrix} 1 & 4 \\ -2 & 6 \end{pmatrix} \Rightarrow c_{31}^A = (-1)^{3+1} \begin{vmatrix} 1 & 4 \\ -2 & 6 \end{vmatrix} = 6 + 8 = 14$$

$$A(1|2) = \begin{pmatrix} -7 & 6 \\ 9 & -3 \end{pmatrix} \Rightarrow c_{12}^A = (-1)^{1+2} \begin{vmatrix} -7 & 6 \\ 9 & -3 \end{vmatrix} = -(21 - 54) = 33$$

\ul{Ví dụ 1}: Cho
$$A = \begin{pmatrix} 2 & -3 & 8 \\ 1 & 6 & -2 \\ -4 & -7 & 5 \end{pmatrix}$$. Tìm
$$\begin{cases} A(3|2) = ? \Rightarrow c_{32}^A = ? \\ A(1|3) = ? \Rightarrow c_{13}^A = ? \\ A(2|1) = ? \Rightarrow c_{21}^A = ? \end{cases}$$

\ul{Ví dụ 2}: Cho
$$A = \begin{pmatrix} -8 & 5 & -1 \\ -6 & 10 & 3 \\ -2 & 4 & -7 \end{pmatrix}$$. Tìm
$$\begin{cases} A(2|3) = ? \Rightarrow c_{23}^A = ? \\ A(3|1) = ? \Rightarrow c_{31}^A = ? \\ A(1|2) = ? \Rightarrow c_{12}^A = ? \end{cases}$$

\ul{Ví dụ 3}: Cho
$$C = \begin{pmatrix} 10 & -4 & -1 \\ -5 & 17 & -2 \\ 1 & -8 & -6 \end{pmatrix}$$. Tìm
$$\begin{cases} C(3|2) = ? \Rightarrow c_{32}^C = ? \\ C(1|3) = ? \Rightarrow c_{13}^C = ? \\ C(2|1) = ? \Rightarrow c_{21}^C = ? \end{cases}$$

\ul{Ví dụ 4}: Cho
$$D = \begin{pmatrix} -2 & 1 & 0 & -6 \\ 8 & -5 & -3 & 2 \\ -10 & 4 & 7 & -9 \end{pmatrix}$$. Tìm
$$\begin{cases} D(1|4) = ? \Rightarrow c_{14}^D = ? \\ D(3|2) = ? \Rightarrow c_{32}^D = ? \end{cases}$$

\ul{Ví dụ 5}: Cho
$$E = \begin{pmatrix} -2 & 5 & -10 & 8 \\ 4 & -3 & 7 & 1 \\ 6 & 0 & -4 & -2 \\ -9 & 1 & 5 & 11 \end{pmatrix}$$. Tìm
$$\begin{cases} E(4|4) = ? \Rightarrow c_{44}^E = ? \\ E(2|3) = ? \Rightarrow c_{23}^E = ? \end{cases}$$

* Tính $\det(A) = |A|$

Ta có thể tính định thức của
$A$ dựa theo bất
kỳ dòng nào hoặc cột nào của
$A$. Cụ thể như
sau:

Xét $A = (a_{ij})_{n \times n}$ thì

$\det(A) = |A| = a_{i1}c_{i1}^A + a_{i2}c_{i2}^A + \dots + a_{in}c_{in}^A$

$= a_{1j}c_{1j}^A + a_{2j}c_{2j}^A + \dots + a_{nj}c_{nj}^A$

\ldots{}

$= a_{i1}c_{i1}^A + a_{i2}c_{i2}^A + \dots + a_{in}c_{in}^A$

$= a_{1j}c_{1j}^A + a_{2j}c_{2j}^A + \dots + a_{nj}c_{nj}^A$

\ul{Ví dụ mẫu}: Cho
$$A = \begin{bmatrix} -3 & 5 & -2 \\ 10 & 6 & -1 \\ 4 & -5 & 8 \end{bmatrix}$$

Ta có: $\det(A) = |A|$ (dòng
3) =

$$4c_{31}^A + (-5)c_{32}^A + 8c_{33}^A = 4.(-1)^{3+1} \begin{vmatrix} 5 & -2 \\ 6 & -1 \end{vmatrix} + (-5).(-1)^{3+2} \begin{vmatrix} -3 & -2 \\ 10 & -1 \end{vmatrix} + 8.(-1)^{3+3} \begin{vmatrix} -3 & 5 \\ 10 & 6 \end{vmatrix}$$

$= 4(-5 + 12) + 5(3 + 20) + 8(-18 - 50) = -401$

\emph{\ul{Lưu ý}}:

+ Ta thường chọn dòng hoặc cột nào có nhiều hệ số = 0 nhất để tính
$\det(A) = |A|$.

+ Nếu ta thấy trong
$A$ có 2 dòng tỷ
lệ nhau (hoặc 2 cột tỷ lệ nhau) thì suy ra
$\det(A) = |A| = 0$.

\end{document}
